\documentclass[11pt,a4paper]{article}

% ------------------------------------------------
% PACKAGES
% ------------------------------------------------
\usepackage[utf8]{inputenc}
\usepackage[T1]{fontenc}
\usepackage[french]{babel}
\usepackage{geometry}
\usepackage{graphicx}
\usepackage{float}
\usepackage{booktabs}
\usepackage{array}
\usepackage{tabularx}
\usepackage{amsmath}
\usepackage{setspace}
\usepackage{titlesec}
\usepackage{fancyhdr}
\usepackage{xcolor}
\usepackage{hyperref}
\usepackage{caption}
\usepackage{enumitem}
\usepackage{tikz}
\usepackage{eso-pic}

% ------------------------------------------------
% PAGE DESIGN
% ------------------------------------------------
\geometry{
    left=2.5cm,
    right=2.5cm,
    top=2.5cm,
    bottom=2.5cm
}

\onehalfspacing

\definecolor{navy}{RGB}{25,50,85}
\definecolor{blue}{RGB}{45,100,155}

\hypersetup{
    colorlinks=true,
    linkcolor=navy,
    citecolor=blue,
    urlcolor=blue
}

\titleformat{\section}
{\Large\bfseries\color{navy}}
{\thesection}{1em}{}

\titleformat{\subsection}
{\large\bfseries\color{blue}}
{\thesubsection}{1em}{}

\pagestyle{fancy}
\fancyhf{}
\fancyhead[L]{Projet Data Science — LPI 2023}
\fancyhead[R]{ENCG Settat}
\fancyfoot[C]{\thepage}
\setlength{\headheight}{14pt}

% ------------------------------------------------
% BORDURE DISCRETE SUR TOUTES LES PAGES
% ------------------------------------------------
\AddToShipoutPictureBG{%
    \begin{tikzpicture}[remember picture,overlay]
        \draw[navy,line width=0.6pt]
        ([xshift=0.9cm,yshift=-0.9cm]current page.north west)
        rectangle
        ([xshift=-0.9cm,yshift=0.9cm]current page.south east);
    \end{tikzpicture}%
}

% ------------------------------------------------
% FIGURE COMMAND
% ------------------------------------------------
\newcommand{\projectfigure}[3]{
\begin{figure}[H]
    \centering
    \includegraphics[width=0.85\textwidth]{figures/#1}
    \caption{#2}
    \label{#3}
\end{figure}
}

% ------------------------------------------------
% DOCUMENT
% ------------------------------------------------
\begin{document}

% =================================================
% PAGE DE GARDE
% =================================================

\begin{titlepage}
\centering

\vspace*{1cm}

{\Large \textbf{ÉCOLE NATIONALE DE COMMERCE ET DE GESTION}}\\[0.3cm]
{\large ENCG SETTAT}\\[1.5cm]

{\color{navy}\rule{\textwidth}{1.5pt}}\\[0.5cm]

{\Huge\bfseries Analyse de la performance}\\[0.2cm]
{\Huge\bfseries logistique des pays}\\[0.5cm]

{\Large À partir du Logistics Performance Index (LPI) 2023}\\[0.5cm]

{\color{navy}\rule{\textwidth}{1.5pt}}\\[1.5cm]

{\Large \textbf{Projet de Data Science}}\\[1cm]

\begin{tabular}{rl}
\textbf{Réalisé par :} & AHNNAOU Youness\\
                       & AIT BOUSSALEM Israa\\[0.5cm]
\textbf{Encadré par :} & Abderrahim Larhlimi\\[0.3cm]
\textbf{Formation :} & Master 1 -- Commerce International\\[0.3cm]
\textbf{Établissement :} & ENCG Settat\\[0.3cm]
\textbf{Année universitaire :} & 2026--2027\\
\end{tabular}

\vfill

{\large Année universitaire 2026--2027}

\end{titlepage}

% =================================================
% RESUME
% =================================================

\pagenumbering{roman}

\section*{Résumé}
\addcontentsline{toc}{section}{Résumé}

Ce projet porte sur l'analyse de la performance logistique des pays
à partir des données du Logistics Performance Index (LPI) 2023
publiées par la Banque mondiale.

L'objectif est d'identifier des groupes de pays présentant des
caractéristiques logistiques similaires en mobilisant des méthodes
de clustering non supervisé.

L'étude repose sur six composantes du LPI : les douanes,
les infrastructures, les expéditions internationales, la qualité
des services logistiques, le suivi et la traçabilité, ainsi que
la ponctualité.

Après les étapes de préparation des données, d'analyse exploratoire
et de standardisation, plusieurs configurations de clustering ont
été étudiées. Le modèle K-Means avec trois clusters a finalement
été retenu pour son intérêt en matière d'interprétation économique.

Les résultats permettent de distinguer trois catégories de pays :
à faible performance logistique, à performance intermédiaire et
à performance élevée.

\vspace{0.5cm}

\noindent
\textbf{Mots-clés :} Data Science, performance logistique,
LPI 2023, clustering, K-Means, ACP, commerce international.

\newpage

\tableofcontents
\newpage

\listoffigures
\newpage

\pagenumbering{arabic}

% =================================================
% INTRODUCTION
% =================================================

\section*{Introduction générale}
\addcontentsline{toc}{section}{Introduction générale}

Dans un contexte marqué par l'intensification des échanges
internationaux, la performance logistique constitue un facteur
important de compétitivité pour les économies nationales.
La capacité d'un pays à assurer le transport, le suivi et la
livraison des marchandises influence directement la fluidité
des échanges commerciaux.

Les différences entre les systèmes logistiques nationaux peuvent
être étudiées à travers le Logistics Performance Index (LPI),
un indicateur publié par la Banque mondiale permettant d'évaluer
plusieurs dimensions de la performance logistique.

Dans ce contexte, notre projet mobilise des méthodes de Data
Science afin d'identifier des groupes de pays présentant des
profils logistiques comparables.

\subsection*{Problématique}

Comment les pays peuvent-ils être regroupés en fonction des six
composantes de leur performance logistique, selon les données
du Logistics Performance Index 2023 ?

\subsection*{Objectifs du projet}

Les principaux objectifs sont :

\begin{itemize}
    \item Comprendre les différentes dimensions de la performance
    logistique.
    \item Explorer et préparer les données du LPI 2023.
    \item Identifier des groupes de pays présentant des profils
    logistiques similaires.
    \item Évaluer les résultats de plusieurs méthodes de clustering.
    \item Interpréter les groupes obtenus dans une perspective
    économique et commerciale.
\end{itemize}

% =================================================
% CHAPITRE 1
% =================================================

\section{Cadre conceptuel : le Logistics Performance Index}

\subsection{Présentation du LPI}

Le Logistics Performance Index est un indicateur développé
par la Banque mondiale afin d'évaluer la performance des systèmes
logistiques des pays.

Il permet notamment d'étudier les conditions dans lesquelles
les marchandises circulent à travers les frontières et au sein
des chaînes logistiques internationales.

\subsection{Les six composantes étudiées}

Notre analyse repose sur les six dimensions suivantes :

\begin{enumerate}
    \item \textbf{Customs :} efficacité des procédures douanières.
    \item \textbf{Infrastructure :} qualité des infrastructures
    de transport et de logistique.
    \item \textbf{International Shipments :} facilité d'organisation
    des expéditions internationales.
    \item \textbf{Logistics Competence and Quality :} qualité et
    compétence des services logistiques.
    \item \textbf{Tracking and Tracing :} capacité de suivi et de
    traçabilité des expéditions.
    \item \textbf{Timeliness :} respect des délais de livraison.
\end{enumerate}

Les scores de ces composantes sont exprimés sur une échelle
allant de 1 à 5.

% =================================================
% CHAPITRE 2
% =================================================

\section{Présentation et préparation des données}

\subsection{Source des données}

Les données utilisées proviennent de la Banque mondiale et
correspondent au Logistics Performance Index 2023.

Le fichier initial contient plusieurs observations et variables
relatives aux pays, aux indicateurs et aux périodes étudiées.

L'analyse a été centrée sur les observations de l'année 2023
et sur les six composantes du LPI.

\subsection{Sélection des observations et des variables}

Après les opérations de filtrage et de transformation,
un tableau structuré a été constitué avec une observation
par pays et six variables quantitatives.

L'échantillon final comprend 139 pays.

Les identifiants des pays ont été conservés pour l'interprétation
des résultats, mais n'ont pas été utilisés comme variables
d'entrée du clustering.

\subsection{Nettoyage et vérification}

Plusieurs contrôles ont été effectués :

\begin{itemize}
    \item Vérification des valeurs manquantes.
    \item Détection des éventuels doublons.
    \item Vérification de la cohérence des scores.
    \item Contrôle de l'échelle de mesure des variables.
\end{itemize}

Aucune valeur manquante n'a été détectée dans les six variables
retenues pour la segmentation.

\subsection{Standardisation}

Les six variables possèdent la même échelle de mesure.
Une standardisation a néanmoins été appliquée afin de centrer
et réduire les variables avant le clustering.

La méthode StandardScaler transforme chaque variable selon :

\[
Z = \frac{X-\mu}{\sigma}
\]

où $\mu$ représente la moyenne et $\sigma$ l'écart-type
de la variable.

Cette transformation permet de rendre les variables
comparables dans le calcul des distances.

% =================================================
% CHAPITRE 3
% =================================================

\section{Analyse exploratoire des données}

L'analyse exploratoire vise à comprendre la distribution des
scores logistiques et les relations entre les six composantes.

\subsection{Distribution des scores}

L'étude des distributions permet d'observer les différences
de niveaux de performance entre les pays.

\projectfigure
{figure_01.png}
{Distribution des scores des six composantes du LPI.}
{fig:distribution}

\projectfigure
{figure_02.png}
{Dispersion des scores des six composantes du LPI à travers les pays (diagrammes en boîte).}
{fig:boxplots}

\subsection{Relations entre les composantes}

L'étude des corrélations permet d'identifier les relations
entre les différentes dimensions de la performance logistique.

\projectfigure
{figure_03.png}
{Matrice de corrélation des six composantes.}
{fig:correlation}

Ces relations permettent de mieux comprendre les liens entre
les infrastructures, les procédures douanières et la qualité
des services logistiques.

\projectfigure
{figure_04.png}
{Scores moyens des six composantes du LPI en 2023.}
{fig:average-scores}

\projectfigure
{figure_05.png}
{Les dix pays présentant les scores moyens les plus élevés sur les six composantes du LPI.}
{fig:top10}

\projectfigure
{figure_06.png}
{Les dix pays présentant les scores moyens les plus faibles sur les six composantes du LPI.}
{fig:bottom10}

% =================================================
% CHAPITRE 4
% =================================================

\section{Méthodologie de clustering}

\subsection{Principe du clustering non supervisé}

Le clustering est une méthode d'apprentissage non supervisé
qui consiste à regrouper les observations présentant des
caractéristiques similaires.

Contrairement à l'apprentissage supervisé, cette approche
ne nécessite pas de variable cible prédéfinie.

Dans notre étude, les pays sont regroupés à partir de leurs
scores sur les six composantes du LPI.

\subsection{Algorithme K-Means}

K-Means est un algorithme de clustering qui cherche à répartir
les observations en $K$ groupes en minimisant la somme des
distances au carré entre les observations et le centre de
leur cluster.

L'objectif est de construire des groupes relativement homogènes
à l'intérieur et distincts les uns des autres.

\projectfigure
{figure_07.png}
{Répartition des pays selon les clusters K-Means, illustrée par les scores en douanes et en infrastructures.}
{fig:clusters-customs-infra}

\subsection{Clustering hiérarchique}

Une méthode de clustering hiérarchique a également été étudiée
afin de comparer les résultats obtenus avec ceux de K-Means.

Cette approche permet de construire progressivement des groupes
en fonction de la proximité entre les observations.

\subsection{Détermination du nombre de clusters}

Plusieurs valeurs de $K$ ont été étudiées afin d'identifier
une segmentation pertinente.

Les méthodes d'évaluation mobilisées comprennent :

\begin{itemize}
    \item La méthode du coude (Elbow Method).
    \item Le coefficient de silhouette.
    \item L'indice de Davies-Bouldin.
    \item L'indice de Calinski-Harabasz.
\end{itemize}

\projectfigure
{figure_10.png}
{Synthèse graphique des critères d'évaluation et de stabilité des configurations K-Means.}
{fig:elbow}

% =================================================
% CHAPITRE 5
% =================================================

\section{Évaluation et choix du modèle final}

La comparaison des configurations étudiées montre que le modèle
à deux clusters présente de meilleurs résultats selon les
métriques statistiques examinées.

Pour $K=2$, le coefficient de silhouette atteint 0,565,
l'indice de Davies-Bouldin est de 0,604 et l'indice de
Calinski-Harabasz est de 317,363.

Pour $K=3$, le coefficient de silhouette est de 0,432,
l'indice de Davies-Bouldin est de 0,780 et l'indice de
Calinski-Harabasz est de 276,246.

Ainsi, la configuration à deux clusters présente une meilleure
séparation statistique selon ces indicateurs.

Toutefois, notre objectif ne se limite pas à obtenir la
meilleure séparation mathématique. Nous souhaitons également
identifier des catégories de performance interprétables
sur le plan économique.

C'est pourquoi nous retenons finalement K-Means avec trois
clusters comme modèle principal d'interprétation.

Ce choix permet de distinguer trois niveaux de performance :
faible, intermédiaire et élevée.

% =================================================
% CHAPITRE 6
% =================================================

\section{Résultats et interprétation économique}

\subsection{Présentation de la segmentation finale}

Le modèle K-Means avec trois clusters a été retenu pour
l'analyse économique.

La répartition obtenue est la suivante :

\begin{table}[H]
\centering
\caption{Répartition des pays selon les trois clusters}
\begin{tabular}{lrr}
\toprule
Catégorie & Nombre de pays & Pourcentage \\
\midrule
Faible performance & 61 & 43,88\% \\
Performance intermédiaire & 41 & 29,50\% \\
Performance élevée & 37 & 26,62\% \\
\midrule
Total & 139 & 100\% \\
\bottomrule
\end{tabular}
\end{table}

\subsection{Visualisation des groupes}

L'Analyse en Composantes Principales (ACP) a été utilisée
pour représenter les pays dans un espace de dimension réduite.

\projectfigure
{figure_08.png}
{Variance expliquée et variance cumulée par les composantes principales.}
{fig:pca-variance}

Elle facilite la visualisation des groupes obtenus tout en
préservant une grande partie de l'information contenue dans
les variables initiales.

\projectfigure
{figure_09.png}
{Projection des pays selon les deux premières composantes principales.}
{fig:pca}


\subsection{Profils logistiques moyens}

Le tableau suivant présente les scores moyens observés pour
chaque catégorie.

\begin{table}[H]
\centering
\caption{Profils moyens des trois catégories}
\small
\begin{tabular}{lrrrrrr}
\toprule
Catégorie & Douanes & Infra. & Expéd. & Qualité & Suivi & Délais \\
\midrule
Faible & 2,264 & 2,284 & 2,464 & 2,446 & 2,428 & 2,746 \\
Intermédiaire & 2,851 & 2,995 & 3,061 & 3,120 & 3,180 & 3,346 \\
Élevée & 3,627 & 3,892 & 3,535 & 3,889 & 3,935 & 3,946 \\
\bottomrule
\end{tabular}
\end{table}

\projectfigure
{figure_11.png}
{Comparaison des scores moyens des composantes par catégorie.}
{fig:profiles}

\projectfigure
{figure_12.png}
{Segmentation finale des pays avec K-Means à trois clusters projetée sur le plan factoriel de l'ACP.}
{fig:pca-final}

\subsection{Interprétation économique}

\textbf{Catégorie à faible performance :}

Cette catégorie regroupe 61 pays. Les scores moyens sont
relativement faibles sur les six dimensions, avec des niveaux
particulièrement bas pour les douanes et les infrastructures.

Ces résultats mettent en évidence des difficultés potentielles
liées à l'efficacité des procédures et à la qualité des
équipements logistiques.

\textbf{Catégorie à performance intermédiaire :}

Cette catégorie comprend 41 pays. Elle présente des scores
supérieurs à ceux du groupe à faible performance, mais reste
en retrait par rapport au groupe le plus performant.

Elle correspond à des systèmes logistiques présentant un
niveau de développement intermédiaire.

\textbf{Catégorie à performance élevée :}

Cette catégorie rassemble 37 pays. Elle se caractérise par
des scores moyens élevés dans les six dimensions étudiées.

Les résultats traduisent une meilleure performance globale
des systèmes logistiques de ces pays.

\subsection{Analyse complémentaire des profils}

La comparaison des profils permet d'identifier les dimensions
sur lesquelles les écarts entre catégories sont les plus
importants.

Cette analyse peut aider à orienter les réflexions sur les
priorités d'amélioration des systèmes logistiques nationaux.

\projectfigure
{figure_13.png}
{Heatmap des profils logistiques standardisés par catégorie.}
{fig:heatmap}

% =================================================
% DISCUSSION
% =================================================

\section{Discussion et limites de l'étude}

Les résultats obtenus mettent en évidence l'existence de
plusieurs profils logistiques parmi les pays étudiés.

La segmentation en trois catégories facilite la lecture
économique des différences entre systèmes logistiques.

Néanmoins, certaines limites doivent être prises en compte.

Premièrement, le clustering dépend du choix des variables
et de la méthode de standardisation.

Deuxièmement, les résultats dépendent du nombre de clusters
retenu. La configuration à deux groupes présente de meilleurs
indicateurs statistiques, tandis que celle à trois groupes
offre une interprétation économique plus détaillée.

Troisièmement, les groupes obtenus décrivent des similarités
statistiques et ne permettent pas, à eux seuls, d'établir
des relations de causalité.

Enfin, l'étude repose sur les données de l'année 2023 et
ne permet pas d'analyser directement l'évolution temporelle
des performances logistiques.

% =================================================
% CONCLUSION
% =================================================

\section*{Conclusion générale}
\addcontentsline{toc}{section}{Conclusion générale}

Ce projet avait pour objectif de regrouper les pays selon
les six composantes de leur performance logistique à partir
des données du Logistics Performance Index 2023.

La démarche adoptée a combiné l'exploration des données,
leur préparation, la standardisation, plusieurs méthodes
de clustering et l'interprétation économique des résultats.

Le modèle K-Means à trois clusters a permis de distinguer
trois catégories : faible performance, performance
intermédiaire et performance élevée.

Bien que la configuration à deux clusters présente de
meilleurs résultats statistiques, le choix de trois clusters
a été privilégié pour sa capacité à fournir une segmentation
plus détaillée et plus exploitable sur le plan économique.

Cette étude montre ainsi l'intérêt des méthodes de Data
Science pour analyser les différences de performance
logistique entre les pays et mieux comprendre les enjeux
associés à la compétitivité du commerce international.

Des recherches futures pourraient intégrer plusieurs années
de données ou d'autres indicateurs économiques afin
d'approfondir cette analyse.

% =================================================
% BIBLIOGRAPHIE
% =================================================

\newpage
\begin{thebibliography}{9}

\bibitem{worldbank}
World Bank.
\textit{Logistics Performance Index (LPI)}.
Banque mondiale.
\url{https://lpi.worldbank.org/}

\bibitem{sklearn}
Scikit-learn.
\textit{K-Means clustering}.
Documentation officielle.
\url{https://scikit-learn.org/stable/modules/clustering.html}

\bibitem{pca}
Scikit-learn.
\textit{Principal Component Analysis (PCA)}.
Documentation officielle.
\url{https://scikit-learn.org/stable/modules/decomposition.html}

\bibitem{pandas}
Pandas.
\textit{Pandas Documentation}.
\url{https://pandas.pydata.org/docs/}

\bibitem{numpy}
NumPy.
\textit{NumPy Documentation}.
\url{https://numpy.org/doc/}

\end{thebibliography}

\end{document}